\subsection{Finite generation of all central supports and the exact zero-Euler class}

The data in the following construction include analytic
neighbourhoods and their transition maps. This is essential
when an inverse operation uses a normal extension or higher
transverse terms. We first specify those data without
referring to a family to be reconstructed.

\paragraph{Finite admissible data.}
The following are the inputs and the checks on them.
\begin{enumerate}
\item An actual proper smooth-total core family
\(f_*:X_*\to\Delta_s\), with its labelled analytic
charts and overlaps and its punctured torus-family
comparison, is given. Its full reduced central space
is one of the cores in UN85: a smooth torus; a
canonical-flat marked Hopf/degree-zero-ruled cycle;
a positive periodic toric log space; or an embedded
rooted-tree or unicyclic elliptic quasi-bundle/Hopf
support of UN61--UN62. The core chart functions,
including their transverse terms, are retained.

\item A finite inverse overlap-compatible sigma record
on that actual core is given. Each arrow includes
the local blowup or contraction map, its centre
and ideal, the target charts, and the isomorphism
on the restored overlap. The maps and their inverse
comparisons are checked on these charts as in UN75.
The resulting space \(Z\) is smooth and proper over
\(\Delta_s\), with full reduced SNC central divisor.
The record is an isomorphism on the punctured family.

\item A finite cyclic birational action on the core,
over \(s\mapsto\zeta s\), is given, where \(\zeta\)
has order \(N\). Conjugating it by the specified
inverse record gives holomorphic maps on every
chart of \(Z\), including the central divisor.
Their compositions satisfy the cyclic group law.
Thus it is an actual regular action of
\(G=\mu_N\) on \(Z\), not only a central action
or an action on its meromorphic function field.

\item A labelled equivariant toric coarsening
\(q:Z\to Y_N\) is given. Its local semigroup
charts, units and maps are explicit normalized
monomial-ideal blowups in the reverse direction.
The finite quotient charts of \(Y_N\) have
invariant ring \(\OO_Y\), with smooth coordinates
in which the descended function is
\[
                  t=u\prod_{i=1}^r x_i^{m_i},
              \qquad r\le3,\quad u\in\OO_Y^*,
              \quad m_i\mid N.
 \tag{C.AR.1}
\]
The normal charts upstairs are exactly the
normalization of \(s^N=t\) in these coordinates.
The labelled height-one subdivisions defining
\(q\) are unimodular, and their overlap maps
and finite action are retained. These ring and
coordinate conditions verify the smooth quotient
\(Y=Y_N/G\) and its full SNC divisor with
the displayed multiplicities.

\item A finite sequence \(\pi:Y\to X\) of
ordinary inverse point or smooth-curve blowups
is given, with their actual smooth target charts
and gluing maps. Their centres are in the central
fibre. The function \(t\) descends through each
map, as is checked by its displayed pullback.
The final total space \(X\) is smooth. Its
central support is allowed arbitrary singularities
and arbitrary full multiplicities.
\end{enumerate}
All quotient, contraction and overlap requirements here
are checks on the supplied finite maps. They do not
refer to the existence of an additional original family.
An ordinary arrow can equivalently be supplied by its
Rees-algebra charts and the corresponding proper map.
For compact data, the core family and base, the cyclic
base cover and finite action, and all maps outside these
central charts are supplied globally; the local overlaps
agree with them on an annulus. This is the compact
version of the same definition.

\begin{statement}{C.2.UN.86}{Proved}{Complete rooted generation and the zero-Euler criterion for arbitrary fibres}
Every proper degeneration of connected complex two-tori
from a smooth threefold, with arbitrary full central
fibre, is the output of the finite admissible construction
above. Conversely every such input produces a proper
degeneration of connected complex two-tori from a smooth
threefold. The inverse record recovers the supplied
family exactly. For compact threefolds over compact
curves, the same assertions hold for the compact data.

Here is an explicit formula selecting precisely the
zero-Euler outputs. Let \(S_*=f_*^{-1}(0)\), let
\(D_i\) be the quotient components in (C.AR.1), and put
\[
 \begin{split}
 D_I^\circ&=\bigcap_{i\in I}D_i
                    \setminus\bigcup_{j\notin I}D_j,\\
 d_I&=\gcd_{i\in I}m_i,\qquad
 \nu_I=\frac{N^{|I|-1}d_I}{\prod_{i\in I}m_i},\\
 R_N&=\sum_I(d_I\nu_I-1)e_c(D_I^\circ).
 \end{split}
 \tag{C.AR.2}
\]
For the ordinary resolution record \(\pi\), define
\[
 \Delta_\pi=2p_\pi+\sum_\alpha(2-2g(C_\alpha)).
 \tag{C.AR.3}
\]
Let \(\Delta_\sigma\) be the signed normal-cone Euler
change in the forward direction \(Z\dashrightarrow X_*\)
of the retained sigma record. Then its output satisfies
\[
 e\bigl(f^{-1}(0)_{\mathrm{red}}\bigr)
       =e(S_*)-\Delta_\sigma-R_N-\Delta_\pi.
 \tag{C.AR.4}
\]
Every zero-Euler fibre is therefore obtained, and only
such fibres are obtained, by the admissible data with
\[
 \Delta_\sigma+R_N+\Delta_\pi=
 \begin{cases}
  2[N_0:\Gamma_0],&\text{for the positive toric core},\\
  0,&\text{for each other core}.
 \end{cases}
 \tag{C.AR.5}
\]
Here \(N_0\) and \(\Gamma_0\) are the two-dimensional
affine lattices of the toric core; the root exponent
\(N\) in (C.AR.1) is a separate integer. The construction
in (C.AR.5) specifies spaces, bundles, finite maps,
centres and analytic gluings, as well as their Euler
numbers. Reversing a retained record recovers every
specified normal extension and every transverse jet.

The choice between a canonical-flat core and a positive
Hopf elliptic-network core is intrinsic to the original
degeneration. Write the original full fibre and the
primitive generator as
\[
        \operatorname{div}(f)=\sum_\alpha n_\alpha S_\alpha,
        \qquad \operatorname{div}(\Omega)
                       =\sum_\alpha a_\alpha S_\alpha.
 \tag{C.AR.6}
\]
Let \(C_\alpha\to T_\alpha\) be the actual finite normal
root cover of the normalization of \(S_\alpha\) defined
by (C.RK.18), and let \(\tau\) be the density threshold
(C.RK.19). Set
\[
 H=\max\left\{\frac{a_\alpha+1}{n_\alpha}:
       C_\alpha\text{ has a component with Hopf minimal model}
       \right\},
 \tag{C.AR.7}
\]
with \(H=\tau\) when this set is empty. Then \(H\ge\tau\).
For every actual divisible semistable root, its highest
Hopf canonical order, when such a component exists, is
\[
                         h=N(H-\tau).
 \tag{C.AR.8}
\]
Thus \(H=\tau\) gives the canonical-flat reduction,
and \(H>\tau\) gives the positive-Hopf reduction of
UN84--UN85. The covers, threshold and this branch test
do not depend on an embedded resolution or on its
chosen divisible root exponent.
\end{statement}

\begin{proof}
\par\medskip\noindent\textbf{The actual forward construction for an arbitrary fibre.}
Shrink the base to a disc about the chosen value.
The proper direct image \(f_*K_X\) is a torsion-free
rank-one sheaf, so it has a primitive generator
\(\Omega\). Its zero divisor is vertical, giving
(C.AR.6). Take an actual embedded resolution
\(\pi:Y\to X\), by smooth point and curve blowups
in the central fibre, with smooth SNC components
\(D_i\). Retain every centre and every blowup map.
The analytic embedded-resolution comparison used
in UN78, from \cite{BierstoneMilman1997}, gives
\[
 \operatorname{div}(f\pi)=\sum_i m_iD_i,
 \qquad \operatorname{div}(\pi^*\Omega)=\sum_i k_iD_i,
 \qquad N=\operatorname{lcm}_i m_i.
 \tag{C.AR.9}
\]

Normalize the literal fibre product for \(t=s^N\)
and use the equivariant semistable maps of UN78:
\[
 Z\xrightarrow{q}Y_N
     =\operatorname{Norm}(Y\times_{\Delta_t}\Delta_s)
                   \longrightarrow Y.
 \tag{C.AR.10}
\]
For clarity, the local root lattice is
\[
 \left\{(v,h)\in\Z^r\oplus\Z:
                      \sum_i m_iv_i=Nh\right\}.
 \tag{C.AR.11}
\]
Its old height-one rays are \(((N/m_i)e_i,1)\).
In dimensions at most three, subdividing each
height-one segment or triangle at its lattice
points gives unimodular cones. These refinements
are implemented by normalized blowups of
monomial ideals, which are unchanged when a
boundary equation is multiplied by a unit.
Their labelled maps therefore glue. The root
group fixes the old rays within a chart and
permutes the charts, so orbitwise choices give
the actual equivariant map (C.AR.10). Retaining
the old component and stratum labels prevents
ray identifications within an old cone. Thus
\(Z\) is smooth and its full fibre is reduced
SNC, including globally embedded components.

The normalized root is connected. Its generic
fibre over \(s\ne0\) is the connected original
torus; all components of the normalized fibre
product dominate the base, so there cannot be
an additional component. Equivalently its
Stein factor has degree one over the normal
disc. All the toric maps are proper and are
isomorphisms away from zero.

Apply the complete reduced theorem UN85 to
\(Z\), and retain its actual comparison
\(\phi:Z\dashrightarrow X_*\). This gives
one of the listed cores and all the inverse
chart maps in the input. The root action
becomes the actual birational action
\[
            \gamma_* =\phi\gamma\phi^{-1}
                            \quad(\gamma\in\mu_N).
 \tag{C.AR.12}
\]
Its group law holds on the dense punctured
family, hence as meromorphic maps. Reversing
\(\phi\) first restores the original regular
action on \(Z\). Applying the retained
equivariant coarsening then gives \(Y_N\).
For a normal local ring \(A\) of \(Y\),
let \(B\) be the integral closure of
\(A[s]/(s^N-f\pi)\). An invariant element
of \(B\) belongs to \(\operatorname{Frac}(A)\)
and is integral over \(A\), so it belongs
to \(A\). The reverse inclusion is immediate.
This remains true when local root branches
split and are permuted by the group. Thus
\(B^{\mu_N}=A\), and the full comparison is
\[
 X_*\dashrightarrow Z\xrightarrow{q}Y_N
              \longrightarrow Y_N/\mu_N=Y
                       \xrightarrow{\pi}X.
 \tag{C.AR.13}
\]
Every arrow and overlap in this diagram has
been retained. In particular the finite
quotient follows the inverse record; no
central quotient of an indeterminate
birational action is taken.

\par\medskip\noindent\textbf{The converse from the specified finite input.}
Begin with the data at the start of the
subsection, independently of any original
\(X\). The inverse sigma maps give the
specified smooth proper \(Z\), and the
chart checks give an actual regular finite
action. Its equivariant coarsening and
invariant-coordinate charts construct the
smooth \(Y\) and its function (C.AR.1).
The subsequent ordinary contraction maps
construct the smooth \(X\) and its descended
function. Properness follows from the
proper core and retained proper maps;
for a local overlap replacement it is the
proper-graph and annular comparison of
UN75. The ordinary target maps and finite
quotient are proper, so the final family
is proper after the common disc has been
shrunk if required.

For a fixed nonzero \(t\), the \(N\)
choices of \(s\) are permuted freely and
transitively by \(\mu_N\). Inclusion of
any one \(Z_s\) into their union induces
a biholomorphism onto the quotient fibre.
The coarsening and ordinary maps are
isomorphisms there. Consequently the
final general fibre is exactly that
connected complex two-torus. Since the
total space is smooth, the nonconstant
base function is a non-zero-divisor;
the map to the smooth curve is flat.
Its Stein factor has generic degree one
and its base is normal, so all its fibres
are connected. These deductions use only
the supplied maps and their checks.

The inverse elementary record is therefore
a constructive description of the resulting
analytic support, with its actual ambient
neighbourhood. Forward construction shows
that every original support has such a
record. The converse shows that every
admissible record produces such a support.
Isomorphic labelled atlases and commuting
record maps give biholomorphic outputs
by their explicit glued isomorphism.
No uniqueness of the reduction record is
required for this exact generation statement.

\par\medskip\noindent\textbf{Original support Euler numbers, including the finite quotient.}
The actual calculation (C.RK.15)--(C.RK.17)
applies to (C.AR.10) without a canonical-order
restriction. Over \(D_I^\circ\), the central
normal root is a finite etale cover of degree
\(d_I\). The toric resolution fibre has Euler
number \(\nu_I\): its zero-dimensional torus
orbits are the unimodular top cones, whose
number is the normalized volume of the
height-one simplex. The exact Bezout
determinant in (C.RK.15) gives the value
in (C.AR.2). Thus
\[
       e(S_Z)-e(S_Y)=R_N.
 \tag{C.AR.14}
\]
The fixed-point average for the actual
finite root group gives the same quotient
comparison, including globally nontrivial
sheet permutations. This is not a free-cover
Euler division at the central fibre.

The ordinary blowups give
\(e(S_Y)-e(S)=\Delta_\pi\), and the
forward sigma record gives
\(e(S_*)-e(S_Z)=\Delta_\sigma\).
For an arrow with centre ideal \(I\),
its signed term is the difference between
the Euler numbers of the actual
\(\operatorname{Proj}(\bigoplus I^a/I^{a+1})\)
and its centre, on the finite constructible
subdivision of UN75. Restored overlaps
cancel. Combining these three exact
comparisons proves (C.AR.4). The core
Euler numbers are those of UN85, proving
(C.AR.5). Since both construction
directions have already been established,
this selects the full class of zero-Euler
outputs, including singular and nonreduced
original fibres.

\par\medskip\noindent\textbf{Intrinsic root data and the two endpoint branches.}
The local finite algebras defining
\(C_\alpha\) are glued from
\(w_j^{n_\alpha}=u_j|_{T_\alpha}\), where
\(f=g_j^{n_\alpha}u_j\); their transition
is \(w_k=(g_j/g_k)w_j\). Their finite
normalization is an actual compact
normal cover. No global meromorphic
frame of a nonalgebraic normal line
is required. For every divisible
root, the strict-cover comparison is
the integral local function
\(w_j=s^{N/n_\alpha}/g_j\), as proved
in (C.RK.21). Every connected component
of this fixed cover is obtained, with
its actual proper birational resolution.

All exceptional surfaces of the original
embedded resolution are projective, and
their finite normal covers are projective
by the coherent-algebra GAGA comparison
in UN78. The additional semistable
components are ruled or rational.
Thus no Hopf component is omitted by
the original covers \(C_\alpha\).

Proper change of variables in the actual
density gives
\[
                   \tau=\min_i\frac{k_i+1}{m_i}.
 \tag{C.AR.15}
\]
At a height-one root ray \((v,1)\),
the pulled canonical form has order
\(\sum_i(k_i+1)v_i-1\). Its minimum
is \(N\tau-1\); the old strict ray
over \(D_i\) has order
\(N(k_i+1)/m_i-1\). Dividing by
\(s^{N\tau-1}\) gives the primitive
root generator, whose strict orders
are therefore
\[
                  N\left(\frac{k_i+1}{m_i}-\tau\right).
 \tag{C.AR.16}
\]
Original strict transforms retain
\((k_i,m_i)=(a_\alpha,n_\alpha)\).
The fixed-cover comparison and exclusion
of exceptional Hopfs now give exactly
(C.AR.8), including its integrality.
When no Hopf occurs, UN75 applies
directly. Otherwise all Hopf root
orders are zero exactly when
\(H=\tau\). If \(H>\tau\), UN84
and UN85 give the positive-Hopf core.
Both the threshold and the covers
were defined on the original family,
so this branch test is independent
of the resolution and exponent.

\par\medskip\noindent\textbf{Compact data and the extent of the core parametrizations.}
For an original compact base \(B\), let
\(c\) be the selected value and let
\(\Sigma\) be the finite critical-value set.
Choose a point \(p\notin\Sigma\cup\{c\}\).
For sufficiently large \(d\),
\[
 H^1\!\left(B,\OO_B(dp-2c-
                    \sum_{v\in\Sigma\setminus\{c\}}v)\right)=0.
\]
Indeed its degree is greater than \(2g(B)-2\),
so Serre duality gives the vanishing. The
actual restriction sequence consequently
surjects from \(H^0(B,\OO_B(dp))\) onto
the prescribed first jet at \(c\) and
values at the other critical points.
Choose the meromorphic function \(a\)
with jet \(a=t+O(t^2)\) in the fixed
coordinate at \(c\), and value one at
every other point of \(\Sigma\).
Its only possible pole is \(p\), and
all its other zeros and poles are regular
values of the original family.

Normalize the compact cyclic base cover
\(s^N=a\). The simple zero at \(c\)
makes it connected and totally ramified
there. It is etale over every other
critical value. Near its additional
branch values the original map is a
submersion, so its base change is locally
a product with the smooth normalized
cover curve and is smooth. Over the
other critical values the base change
is etale over the original smooth total
space. The normalized total space is
therefore smooth away from the selected
root fibre. Apply the preceding root
resolution and modification maps there,
retaining the actual identity comparisons
on its annulus and on the complement.
The local function \(a\) is a coordinate
at \(c\); its exact unit comparison with
\(t\) is retained in these root charts.

The finite action is conjugated by that
actual global modification and becomes
regular again upon reversal. The quotient
and original blowdowns recover the
original compact family. Conversely the
globally supplied compact admissible
data give a compact output by the same
proper maps.

The central support parameters of the
flat cycles and elliptic networks are
the complete marked line, torsor,
finite-character and gluing data in
UN57 and UN61--UN63. The positive
toric core is parametrized by its
periodic affine triangulation and
normalized log gluing data in UN71.
Their compact occurrence is proved
in those statements and UN73.
Those constructions supply actual
compact starting families for records
defined on their neighbourhoods.
For recovery of a different given
family, its actual endpoint atlas
and higher transverse terms remain
in the input. The proof uses its
specified overlap isomorphisms and
does not transfer a contraction
through an equality of central log
data alone. This states exactly
which parameters are classified
centrally and which analytic data
the reversible reconstruction retains.
\end{proof}
